\documentclass[10pt,a4paper]{amsart}
\usepackage{amsmath,amssymb,amsthm,graphicx}
% Retained preamble for the Stacks Project chapter snapshot in this directory.
% The chapter's own preamble.tex is a copy of the shared Stacks Project preamble; the block
% below is sliced MECHANICALLY, line for line, from it. Nothing was added, removed or
% redefined; the wrapper class, the xr-hyper/hyperref wiring and the externaldocument list
% are omitted because the assembler supplies the document class.
%
% lines taken from preamble.tex: 3, 5, 9, 11, 12, 156, 157, 158, ... (51 lines total)
%
% package coverage restored from the chapter's own preamble (2026-10-13): xy, xy-2cell

\usepackage[all]{xy}
\xyoption{2cell}
\UseAllTwocells

\newenvironment{reference}{\comment}{\endcomment}
\newenvironment{slogan}{\comment}{\endcomment}
\newenvironment{history}{\comment}{\endcomment}
\theoremstyle{plain}
\newtheorem{theorem}[subsection]{Theorem}
\newtheorem{proposition}[subsection]{Proposition}
\newtheorem{lemma}[subsection]{Lemma}
\theoremstyle{definition}
\newtheorem{definition}[subsection]{Definition}
\newtheorem{example}[subsection]{Example}
\newtheorem{exercise}[subsection]{Exercise}
\newtheorem{situation}[subsection]{Situation}
\theoremstyle{remark}
\newtheorem{remark}[subsection]{Remark}
\newtheorem{remarks}[subsection]{Remarks}
\numberwithin{equation}{subsection}
\def\lim{\mathop{\mathrm{lim}}\nolimits}
\def\colim{\mathop{\mathrm{colim}}\nolimits}
\def\Spec{\mathop{\mathrm{Spec}}}
\def\Hom{\mathop{\mathrm{Hom}}\nolimits}
\def\Ext{\mathop{\mathrm{Ext}}\nolimits}
\def\SheafHom{\mathop{\mathcal{H}\!\mathit{om}}\nolimits}
\def\SheafExt{\mathop{\mathcal{E}\!\mathit{xt}}\nolimits}
\def\Sch{\mathit{Sch}}
\def\Mor{\mathop{\mathrm{Mor}}\nolimits}
\def\Ob{\mathop{\mathrm{Ob}}\nolimits}
\def\Sh{\mathop{\mathit{Sh}}\nolimits}
\def\NL{\mathop{N\!L}\nolimits}
\def\CH{\mathop{\mathrm{CH}}\nolimits}
\def\proetale{{pro\text{-}\acute{e}tale}}
\def\etale{{\acute{e}tale}}
\def\QCoh{\mathit{QCoh}}
\def\Ker{\mathop{\mathrm{Ker}}}
\def\Im{\mathop{\mathrm{Im}}}
\def\Coker{\mathop{\mathrm{Coker}}}
\def\Coim{\mathop{\mathrm{Coim}}}
\DeclareMathSymbol{\boxtimes}{\mathbin}{AMSa}{"02}
\def\QCohstack{\mathcal{QC}\!\mathit{oh}}
\def\Cohstack{\mathcal{C}\!\mathit{oh}}
\def\Spacesstack{\mathcal{S}\!\mathit{paces}}
\def\Quotfunctor{\mathrm{Quot}}
\def\Hilbfunctor{\mathrm{Hilb}}
\def\Curvesstack{\mathcal{C}\!\mathit{urves}}
\def\Polarizedstack{\mathcal{P}\!\mathit{olarized}}
\def\Complexesstack{\mathcal{C}\!\mathit{omplexes}}
\def\Pic{\mathop{\mathrm{Pic}}\nolimits}
\def\Picardstack{\mathcal{P}\!\mathit{ic}}
\def\Picardfunctor{\mathrm{Pic}}
\def\Deformationcategory{\mathcal{D}\!\mathit{ef}}
\usepackage{cleveref}
\crefname{originallemma}{original lemma}{original lemmas}
\Crefname{originallemma}{Original lemma}{Original lemmas}
\makeatletter
\newlabel{lemma-dual-obstruction}{{07YY}{source}{}{sourcepointer-1}{}}
\newlabel{lemma-dual-obstruction@cref}{{[originallemma][0][]07YY}{[1][0][]source}{}{}{}}
\makeatother
\crefname{originallemma}{original lemma}{original lemmas}
\Crefname{originallemma}{Original lemma}{Original lemmas}
\makeatletter
\newlabel{lemma-openness}{{07YN}{source}{}{sourcepointer-2}{}}
\newlabel{lemma-openness@cref}{{[originallemma][0][]07YN}{[1][0][]source}{}{}{}}
\makeatother
\crefname{originallemma}{original lemma}{original lemmas}
\Crefname{originallemma}{Original lemma}{Original lemmas}
\makeatletter
\newlabel{more-algebra-lemma-dual-perfect-complex}{{07VI}{source}{}{sourcepointer-3}{}}
\newlabel{more-algebra-lemma-dual-perfect-complex@cref}{{[originallemma][0][]07VI}{[1][0][]source}{}{}{}}
\makeatother
\crefname{originallemma}{original lemma}{original lemmas}
\Crefname{originallemma}{Original lemma}{Original lemmas}
\makeatletter
\newlabel{situation-dual}{{07YX}{source}{}{sourcepointer-4}{}}
\newlabel{situation-dual@cref}{{[originallemma][0][]07YX}{[1][0][]source}{}{}{}}
\makeatother
\title{Openness of versality from a perfect dual obstruction theory}
\author{The Stacks Project Authors}
\date{stacks-project:artin.tex:local-snapshot-20261012; GNU FDL 1.2 or later}
\begin{document}
\maketitle
\noindent\textit{Editorial note.} Statement span is lines 4714-4722 of sources/stacks-artin/pinned-original-source.tex; proof span is lines 4724-4749, the complete author proof (24 body lines).\par\medskip
\noindent\textit{Editorial note.} Navigation stubs: 4 cited-but-not-included labels are supplied as non-load-bearing pointers with their pinned tags where the snapshot tag index has them.\par\medskip
\noindent\textit{Editorial note.} extra\_packages is empty on purpose: the retained preamble already loads xy with options and a second declaration clashes in the assembler wrapper.\par\medskip
\noindent\textit{Editorial note.} Text is retained verbatim, including the author's own typos and placeholders.\par\medskip
\noindent\textit{Editorial note.} Queue index 31 of queue\_artin.json (pb=1322, stmt 4714-4722, proof 4724-4749).\par\medskip
\noindent\textit{Editorial note.} stacks-artin wave ar01 (2026-10-14): fresh adjudication accept per lane-s3/artin\_judgments.md; pool per Chief 10:5x direct dispatch (idle-first over S12's 01:5x note, SMC precedent); artin.tex 99 rows/67 with proof judged 46 accept / 21 reject / 5 AGP-overlap avoidance rows; patch-sd-envs contract spans scanned.\par\medskip
\noindent\textit{Allowed standard prerequisites (editorial).}\begin{itemize}
\item Explicit prerequisite boundary. The cited inputs are prerequisites; the argument in the retained span is not replaced by any cited result.
\end{itemize}\medskip
\ifdefined\hypertarget\hypertarget{sourcepointer-1}{}\fi
\noindent\textit{Source reference.} Original lemma 07YY: sources/stacks-artin/pinned-original-source.tex, source lines 4616--4628. Cited by the retained statement or proof but not part of the collected spans; supplied only as a navigation pointer. The argument does not depend on its content.\par\medskip
\ifdefined\hypertarget\hypertarget{sourcepointer-2}{}\fi
\noindent\textit{Source reference.} Original lemma 07YN: sources/stacks-artin/pinned-original-source.tex, source lines 4364--4381. Cited by the retained statement or proof but not part of the collected spans; supplied only as a navigation pointer. The argument does not depend on its content.\par\medskip
\ifdefined\hypertarget\hypertarget{sourcepointer-3}{}\fi
\noindent\textit{Source reference.} Original lemma 07VI: sources/stacks-artin/pinned-original-source.tex, source lines 4749--4749. Cited by the retained statement or proof but not part of the collected spans; supplied only as a navigation pointer. The argument does not depend on its content.\par\medskip
\ifdefined\hypertarget\hypertarget{sourcepointer-4}{}\fi
\noindent\textit{Source reference.} Original lemma 07YX: sources/stacks-artin/pinned-original-source.tex, source lines 4570--4585. Cited by the retained statement or proof but not part of the collected spans; supplied only as a navigation pointer. The argument does not depend on its content.\par\medskip
\begin{lemma}
\label{lemma-dual-openness}
In Situation \ref{situation-dual} assume that (iv) of
Lemma \ref{lemma-dual-obstruction} holds and that $K^\bullet$ is a
perfect object of $D(A)$. In this case, if $x$ is versal at a closed
point $u_0 \in U$ then there exists an open neighbourhood
$u_0 \in U' \subset U$ such that $x$ is versal at every finite type
point of $U'$.
\end{lemma}

\begin{proof}
We may assume that $K^\bullet$ is a finite complex of finite projective
$A$-modules. Thus the derived tensor product with $K^\bullet$ is the
same as simply tensoring with $K^\bullet$. Let
$E^\bullet$ be the dual perfect complex to $K^\bullet$, see
More on Algebra, Lemma \ref{more-algebra-lemma-dual-perfect-complex}.
(So $E^n = \Hom_A(K^{-n}, A)$ with differentials the transpose of the
differentials of $K^\bullet$.) Let $E \in D^{-}(A)$ denote the
object represented by the complex $E^\bullet[-1]$.
Let $\xi \in H^1(\text{Tot}(K^\bullet \otimes_A \NL_{A/\Lambda}))$
be the element constructed in Lemma \ref{lemma-dual-obstruction}
and denote $\xi : E = E^\bullet[-1] \to \NL_{A/\Lambda}$ the corresponding
map (loc.cit.). We claim that the pair $(E, \xi)$ satisfies all the
assumptions of Lemma \ref{lemma-openness} which finishes the proof.

\medskip\noindent
Namely, assumption (i) of Lemma \ref{lemma-openness} follows from conclusion
(1) of Lemma \ref{lemma-dual-obstruction}
and the fact that $H^2(K^\bullet \otimes_A^\mathbf{L} -) =
\Ext^1(E, -)$ by loc.cit. Assumption (ii) of
Lemma \ref{lemma-openness} follows from conclusion (2) of
Lemma \ref{lemma-dual-obstruction}
and the fact that $H^1(K^\bullet \otimes_A^\mathbf{L} -) =
\Ext^0(E, -)$ by loc.cit. Assumption (iii) of Lemma \ref{lemma-openness}
is clear.
\end{proof}

\end{document}
